\chapter{Conclusion} \label{chap:ch9}
\paragraph{}


This document established the foundations and technical direction for addressing centralized multi-robot coordination under the MAPF formulation. Through the analysis of MAPF fundamentals and state-of-the-art coordination paradigms, it became clear that centralized approaches remain particularly well suited for scenarios requiring strong safety, predictability, and global consistency guarantees, despite their inherent scalability challenges.

The detailed study of centralized planning methods highlighted the advantages of TEA* in practical fleet coordination systems, namely its explicit handling of temporal constraints and its suitability for online replanning. However, the sequential and prioritized nature of TEA* was identified as a key limitation, as planning outcomes are highly sensitive to agent ordering and no single prioritization heuristic performs consistently well across different environments and task configurations.

Motivated by this observation, a preliminary solution was proposed that extends the existing TEA*-based planner through parallel exploration of multiple priority orderings. By executing several sequential TEA* instances concurrently and selecting the best resulting plan, the proposed approach aims to improve robustness and solution quality without altering the core architecture of the centralized coordination framework. The discussion of alternative execution and synchronization strategies further emphasized the potential to balance planning latency and exploration depth through tunable parameters.

Overall, the conclusions drawn in this document motivate the dissertation phase, which will focus on validating whether parallel prioritization can effectively mitigate ordering sensitivity in TEA*-based fleet coordination systems under realistic conditions.